\chapter{Observador acoplado, autoinercia y conservaci\'on geneal\'ogica de la informaci\'on}
\label{chap:observador-autoinercia-landauer}

\section{Estado conjunto y residencia del observador}

La descripci\'on din\'amica elemental comprende la secci\'on observada, el
observador realizado y el registro que ambos producen. Escribimos
\begin{equation}
 \boxed{\Gamma=(x,\mathcal O,m)\in
 \mathcal X_{\rm TPK}\times\mathfrak O\times\mathfrak M.}
 \label{eq:rm-obs-joint-state}
\end{equation}
La coordenada \(x\) conserva cilindro, ruta, hoja, orientaci\'on, acarreo,
frontera y posibilidades de prolongaci\'on; \(\mathcal O\) declara el
subsistema lector y su resoluci\'on; \(m\) conserva la memoria distinguible
del acoplamiento. El estado conjunto constituye el dominio primario de la
medici\'on y de la respuesta inercial.

Un observador a profundidad \(n\) se representa mediante una carta
\begin{equation}
 q_{\mathcal O,n}:\mathcal X_{\rm TPK}^{(n)}
 \longrightarrow Y_{\mathcal O,n}.
 \label{eq:rm-obs-observer-chart}
\end{equation}
Para \(x\in\mathcal X_{\rm TPK}^{(n)}\), su horizonte observacional es la
fibra
\begin{equation}
 \mathfrak H_{\mathcal O,n}(x)
 =q_{\mathcal O,n}^{-1}(q_{\mathcal O,n}(x)).
 \label{eq:rm-obs-horizon-fiber}
\end{equation}
La fibra re\'une las secciones indistinguibles para la carta declarada. El
refinamiento conserva la naturalidad cuando existen mapas
\(r_{n+1,n}\) y \(s_{n+1,n}\) tales que
\begin{equation}
 q_{\mathcal O,n}\circ r_{n+1,n}
 =s_{n+1,n}\circ q_{\mathcal O,n+1}.
 \label{eq:rm-obs-observer-naturality}
\end{equation}
Esta identidad hace compatible la observaci\'on finita con la genealog\'ia
profunda: el resultado de truncar una lectura refinada coincide con la
lectura del estado truncado.

\begin{definition}[Observador realizado]
Un observador realizado es una tupla
\begin{equation}
 \mathcal O_a=(m_0,U_a,q_a,\mathcal C_a)
 \label{eq:rm-obs-realized-observer}
\end{equation}
formada por una preparaci\'on \(m_0\), un acoplamiento \(U_a\), un lector
de puntero \(q_a\) y una condici\'on de compatibilidad y prolongaci\'on
\(\mathcal C_a\). Su funci\'on consiste en producir un registro distinguible
y emplearlo como condici\'on de continuaci\'on de la ruta.
\end{definition}

La definici\'on sit\'ua el marco de lectura dentro del sistema. La
coordenada publicada depende del acoplamiento, de la resoluci\'on y de la
memoria; la secci\'on enriquecida conserva la totalidad de la que procede esa
coordenada.

\section{Contexto de medici\'on y condicionamiento por fibras}

Sean \((\Omega_S,\Sigma_S)\) y \((\Omega_M,\Sigma_M)\) espacios medibles
de secciones del sistema y del aparato. El contexto \(a\) emplea un
acoplamiento medible
\begin{equation}
 U_a:\Omega_S\times\Omega_M
 \longrightarrow\Omega_S\times\Omega_M
 \label{eq:rm-obs-measurable-coupling}
\end{equation}
y un lector \(q_a:\Omega_M\to Y_a\). La salida es
\begin{equation}
 \operatorname{Out}_a(x)
 =q_a\!\left(\operatorname{pr}_M U_a(x,m_0)\right).
 \label{eq:rm-obs-output}
\end{equation}
Para un resultado \(y\) se define la fibra medible
\begin{equation}
 \Omega_{a,y}
 =\{x\in\Omega_a:\operatorname{Out}_a(x)=y\}.
 \label{eq:rm-obs-outcome-fiber}
\end{equation}
La actualizaci\'on geneal\'ogica es la restricci\'on
\begin{equation}
 \Omega_a\longmapsto\Omega_{a,y}.
 \label{eq:rm-obs-genealogical-update}
\end{equation}
En profundidad finita, la operaci\'on conserva los prefijos que admiten una
extensi\'on en \(\Omega_{a,y}\). La memoria del resultado se incorpora a la
condici\'on de prolongaci\'on; la ruta recorrida permanece registrada.

Si \(\mu_a(\Omega_{a,y})>0\), la medida condicionada es
\begin{equation}
 \mu_{a,y}(B)
 =\frac{\mu_a(B\cap\Omega_{a,y})}
 {\mu_a(\Omega_{a,y})}.
 \label{eq:rm-obs-conditional-measure}
\end{equation}
Esta expresi\'on constituye la carta probabil\'istica del estado relativo.
Su realizaci\'on espectral se obtiene mediante proyecciones cil\'indricas.

\section{PVM cil\'indrica y correspondencia de L\"uders}

Supongamos que \(\Omega_{a,y}\) es una uni\'on finita de cilindros
disjuntos de una profundidad com\'un \(n\). Sea \(P_y\) la suma de sus
proyectores ortogonales en la realizaci\'on de Hilbert. Para un estado
densidad \(\rho\) y una PVM com\'un de sucesos cil\'indricos, la probabilidad
de Born es
\begin{equation}
 p(y)=\Tr(\rho P_y).
 \label{eq:rm-obs-born}
\end{equation}
Cuando \(p(y)>0\), la actualizaci\'on de L\"uders \cite{Luders1951} es
\begin{equation}
 \rho_y=\frac{P_y\rho P_y}{\Tr(\rho P_y)}.
 \label{eq:rm-obs-luders}
\end{equation}

\begin{theorem}[Equivalencia entre restricci\'on de fibra y L\"uders]
\label{thm:rm-obs-fiber-luders}
Sea \(\mathcal A_n\) el \`algebra de sucesos cil\'indricos a profundidad
\(n\). Supongamos que
\begin{equation}
 \mu_a(B)=\Tr(\rho P_B),
 \qquad B\in\mathcal A_n,
 \label{eq:rm-obs-measure-pvm}
\end{equation}
y que \(P_B\) y \(P_y\) pertenecen a la misma medida espectral proyectiva.
Entonces
\begin{equation}
 \boxed{
 \Tr(\rho_yP_B)
 =\frac{\mu_a(B\cap\Omega_{a,y})}
 {\mu_a(\Omega_{a,y})}.}
 \label{eq:rm-obs-luders-fiber-equivalence}
\end{equation}
\end{theorem}

\begin{proof}
La pertenencia a una PVM com\'un da
\(P_BP_y=P_{B\cap\Omega_{a,y}}\). La ciclicidad de la traza y
\cref{eq:rm-obs-measure-pvm} producen
\[
 \Tr(\rho_yP_B)
 =\frac{\Tr(\rho P_yP_BP_y)}{\Tr(\rho P_y)}
 =\frac{\Tr(\rho P_{B\cap\Omega_{a,y}})}
 {\Tr(\rho P_y)},
\]
que coincide con \cref{eq:rm-obs-conditional-measure}.
\end{proof}

La hip\'otesis de PVM com\'un es sustantiva. Dos contextos mutuamente
insesgados admiten descripciones proyectivas propias y requieren un POVM
ruidoso o una dilataci\'on para una lectura conjunta. La correspondencia de
\cref{thm:rm-obs-fiber-luders} se afirma dentro del contexto espectral
declarado.

\section{Instrumento cu\'antico, dilataci\'on y conservaci\'on del canal}

Sea \(\widehat U_a\) la realizaci\'on unitaria del acoplamiento entre
sistema y aparato, y sea \(|m_0\rangle\) la preparaci\'on del puntero. Una
base refinada \(|m_{y,\alpha}\rangle\) induce los operadores de Kraus
\begin{equation}
 M_{y,\alpha}
 =\langle m_{y,\alpha}|\widehat U_a|m_0\rangle,
 \qquad
 \sum_{y,\alpha}M_{y,\alpha}^{\dagger}M_{y,\alpha}=I.
 \label{eq:rm-obs-kraus}
\end{equation}
El instrumento y su efecto son
\begin{equation}
 \mathcal I_y(\rho)
 =\sum_\alpha M_{y,\alpha}\rho M_{y,\alpha}^{\dagger},
 \qquad
 E_y=\sum_\alpha M_{y,\alpha}^{\dagger}M_{y,\alpha}.
 \label{eq:rm-obs-instrument}
\end{equation}
La probabilidad y el estado relativo toman la forma
\begin{equation}
 p(y)=\Tr(\rho E_y),
 \qquad
 \rho_y=\frac{\mathcal I_y(\rho)}{p(y)}.
 \label{eq:rm-obs-general-conditioned-state}
\end{equation}
El canal no selectivo es
\begin{equation}
 \mathcal E_a(\rho)=\sum_y\mathcal I_y(\rho).
 \label{eq:rm-obs-unconditioned-channel}
\end{equation}

\begin{proposition}[Conservaci\'on exacta del canal global]
\label{prop:rm-obs-global-channel}
La isometr\'ia
\begin{equation}
 V_a:\mathcal H_S\longrightarrow\mathcal H_S\otimes\mathcal H_M,
 \qquad
 V_a|\psi\rangle=\widehat U_a(|\psi\rangle\otimes|m_0\rangle)
 \label{eq:rm-obs-stinespring-isometry}
\end{equation}
satisface \(V_a^\dagger V_a=I\). En consecuencia,
\begin{equation}
 \sum_y p(y)=1,
 \qquad
 \Tr\mathcal E_a(\rho)=\Tr\rho,
 \label{eq:rm-obs-trace-preservation}
\end{equation}
y \(\mathcal E_a\) es completamente positivo y preservador de la traza.
Para un estado puro inicial, el estado conjunto
\(V_a|\psi\rangle\) sigue siendo puro y normalizado.
\end{proposition}

\begin{proof}
La unitariedad de \(\widehat U_a\) y la normalizaci\'on de
\(|m_0\rangle\) dan \(V_a^\dagger V_a=I\). La identidad de completitud de
\cref{eq:rm-obs-kraus} preserva la traza. La forma de Kraus demuestra
positividad completa, y una isometr\'ia env\'ia vectores unitarios en
vectores unitarios.
\end{proof}

La conservaci\'on del canal global significa que la evoluci\'on dilatada
retiene todas las amplitudes y todos los registros. El condicionamiento
selecciona una fibra relativa para el lector que obtuvo \(y\). La suma de
los instrumentos reproduce el canal completo. Esta triple distinci\'on
separa evoluci\'on, lectura y estado relativo.

El teorema de Stinespring \cite{Stinespring1955} garantiza una dilataci\'on isom\'etrica para todo
canal completamente positivo. La identificaci\'on del espacio auxiliar con
la memoria TPK se realiza mediante un entrelazador que conserve ruta,
acarreo, registro y refinamiento. La existencia abstracta de la dilataci\'on
y la genealog\'ia concreta del registro son dos resultados coordinables y
tipados.

\section{Estado relativo y formulaci\'on everettiana}

La formulaci\'on de estado relativo proporciona el antecedente conceptual
de esta construcci\'on \cite{Everett1957}.

En una premedici\'on ideal, sup\'ongase que los estados de puntero
\(\{|m_r\rangle\}_r\) son ortonormales y que \(\widehat U_a\) es una
isometr\'ia sobre el subespacio generado por
\(\{|r\rangle\otimes|m_0\rangle\}_r\). Entonces
\begin{equation}
 \widehat U_a(|r\rangle\otimes|m_0\rangle)
 =|r\rangle\otimes|m_r\rangle.
 \label{eq:rm-obs-ideal-premeasurement}
\end{equation}
Para \(|\psi\rangle=\sum_r c_r|r\rangle\), la linealidad produce
\begin{equation}
 |\Psi'\rangle
 =\sum_r c_r|r\rangle\otimes|m_r\rangle.
 \label{eq:rm-obs-entangled-record}
\end{equation}
El registro \(m_r\) define la secci\'on relativa
\begin{equation}
 \rho_r
 =\frac{Q_r|\Psi'\rangle\langle\Psi'|Q_r}
 {\Tr(|\Psi'\rangle\langle\Psi'|Q_r)},
 \qquad
 Q_r=I_S\otimes|m_r\rangle\langle m_r|.
 \label{eq:rm-obs-relative-state}
\end{equation}
En el espacio geneal\'ogico, el mismo resultado corresponde a
\begin{equation}
 \Omega_{a,r}
 =\{x:\operatorname{Out}_a(x)=r\}.
 \label{eq:rm-obs-relative-fiber}
\end{equation}

\begin{theorem}[Correspondencia de estado relativo]
\label{thm:rm-obs-everett-correspondence}
Bajo las hip\'otesis proyectivas de
\cref{thm:rm-obs-fiber-luders}, el estado relativo de
\cref{eq:rm-obs-relative-state} y la medida condicionada sobre
\(\Omega_{a,r}\) asignan las mismas probabilidades a todos los sucesos de
la PVM cil\'indrica declarada. La isometr\'ia global conserva la superposici\'on
correlacionada de todos los registros.
\end{theorem}

\begin{proof}
La primera afirmaci\'on es
\cref{eq:rm-obs-luders-fiber-equivalence}. La segunda se sigue de
\(V_a^\dagger V_a=I\) y de la expresi\'on lineal
\cref{eq:rm-obs-entangled-record}.
\end{proof}

La correspondencia everettiana alcanza el nivel matem\'atico de los estados
relativos: el estado conjunto conserva los canales y cada registro determina
una fibra condicionada. HMT a\~nade a cada fibra una genealog\'ia compuesta
por prefijos, acarreo, orientaci\'on, frontera y prolongaciones compatibles.
La estabilidad macrosc\'opica de los registros pertenece a la teor\'ia de
decoherencia del aparato; la ontolog\'ia de ramas f\'isicamente aut\'onomas
pertenece a una interpretaci\'on adicional. El teorema anterior conserva su
fuerza con independencia de esa elecci\'on ontol\'ogica.

